%% toc: P2 20 20.7,20.8
\section{Elementary extensions by inclusion, nonstandard models, and countability inside a model}
\label{sec:p02-w02-p02-033-inclusion-and-skolem}

\Cref{thm:p02-upward-ls-elementary-extension} produces an elementary
embedding with proper range, and \cref{def:p02-skolem-paradox} states that a
countable model of set theory can contain a set it regards as uncountable;
\cref{thm:p02-nonstandard-arithmetic-model} is the one nonstandard model
obtained from them.  This section supplies the arguments those statements
leave implicit and a set of compactness constructions of nonstandard models.
The first part passes from an elementary embedding to an elementary extension
in which the original structure is an actual substructure, says exactly where
the embedding form and the inclusion form differ, and shows that none of the
cardinal and size hypotheses of upward L\"owenheim--Skolem can be dropped.
The second part gives a procedure for building models by compactness and
applies it to three theories: true arithmetic with an element different from
every numeral, true arithmetic with a nonzero element divisible by every
prime, and the elementary diagram of the real ordered field with an element
above every real.  For each, the finite satisfiability of the theory is
verified with an explicit witness, and the properties of the resulting
model are proved; the limits of the method are recorded.
The third part proves that every model of a modest set theory satisfies the
existence of an uncountable set because the theory proves it, and that in a
countable model the members of that set form an externally countable
collection whose enumerations are not elements of the model.

The inputs are the inherited first-order results: compactness
(\cref{thm:fol-compactness}), soundness and completeness
(\cref{thm:fol-soundness,thm:fol-completeness}), downward
L\"owenheim--Skolem (\cref{thm:downward-ls}), invariance of satisfaction
under isomorphism (\cref{lem:iso-invariance}), the Tarski--Vaught test
(\cref{thm:tarski-vaught}), the diagram tests (\cref{thm:diagram-tests}),
Cantor's theorem in \cref{sec:set-prerequisites-05}, the language and
standard structure of arithmetic (\cref{def:language-arithmetic}), the
compactness construction of nonstandard models of arithmetic
(\cref{thm:nonstandard-model}), and the axioms of
\cref{sec:zf-zfc-01}; their source intervals are
\cite[lines 533--602, 3519--3680, 3927--4061, 4109--4213, 4713--4967,
5043--5131, 6070--6421, 8306--8438, 9188--9261 and
10089--10271]{FrozenFormalLogicCompactnessApplications}.
The arguments below are editorial completions built from those results; they
are not attributed to any external source.
Throughout, homomorphisms, embeddings, isomorphisms, elementary maps,
substructures \(\mathcal M\subseteq\mathcal N\) and elementary substructures
\(\mathcal M\preccurlyeq\mathcal N\) are as defined in
\cref{sec:maps-elementarity-01,sec:maps-elementarity-02,sec:maps-elementarity-03}.

\subsection{From an elementary embedding to an elementary extension}
\label{sec:p02-w02-p02-033-inclusion}

\begin{hbdefinition}{def:p02-w02-p02-033-two-forms}{Embedding form and inclusion form of an elementary extension}
Let \(\mathcal M\) be an \(\Lang\)-structure and \(\kappa\) a cardinal.
\begin{enumerate}
  \item An \emph{elementary extension of \(\mathcal M\) in embedding form} of
        cardinality \(\kappa\) is a pair \((\mathcal N,j)\) of an
        \(\Lang\)-structure \(\mathcal N\) with
        \(\abs{\dom\mathcal N}=\kappa\) and an elementary embedding
        \(j:\mathcal M\to\mathcal N\).  It is \emph{proper} when
        \(j[\dom\mathcal M]\ne\dom\mathcal N\).
  \item An \emph{elementary extension of \(\mathcal M\) in inclusion form}
        of cardinality \(\kappa\) is an \(\Lang\)-structure \(\mathcal N\)
        with \(\abs{\dom\mathcal N}=\kappa\) and
        \(\mathcal M\preccurlyeq\mathcal N\); that is,
        \(\dom\mathcal M\subseteq\dom\mathcal N\), \(\mathcal M\) is a
        substructure of \(\mathcal N\), and the inclusion map is elementary.
        It is \emph{proper} when \(\dom\mathcal M\ne\dom\mathcal N\).
\end{enumerate}
An inclusion-form extension yields an embedding-form extension, with the
inclusion map as \(j\) and the same properness.  The converse passage is
\cref{thm:p02-w02-p02-033-copy-over-image}.
\end{hbdefinition}

\begin{hblemma}{thm:p02-w02-p02-033-elementary-maps-embed}{proved=annex:p02-w02-p02-033-embed}{Elementary embeddings are embeddings}
Let \(h:\mathcal M\to\mathcal N\) be an elementary embedding of
\(\Lang\)-structures.  Then \(h\) is an embedding, and every
\(\Lang\)-sentence true in \(\mathcal M\) is true in \(\mathcal N\).  If
moreover \(\dom\mathcal M\subseteq\dom\mathcal N\) and \(h\) is the
inclusion map, then \(\mathcal M\subseteq\mathcal N\) and
\(\mathcal M\preccurlyeq\mathcal N\).
\end{hblemma}

\begin{hbtheorem}{thm:p02-w02-p02-033-copy-over-image}{proved=annex:p02-w02-p02-033-copy}{Isomorphic copy over the image of an elementary embedding}
Let \(j:\mathcal M\to\mathcal N\) be an elementary embedding of
\(\Lang\)-structures, and put
\(D=\dom\mathcal N\setminus j[\dom\mathcal M]\).
\begin{enumerate}
  \item There are an \(\Lang\)-structure \(\mathcal N'\) and an isomorphism
        \(\pi:\mathcal N'\to\mathcal N\) such that
        \(\dom\mathcal M\subseteq\dom\mathcal N'\),
        \(\pi(a)=j(a)\) for every \(a\in\dom\mathcal M\), and
        \(\mathcal M\preccurlyeq\mathcal N'\).
  \item For any such \(\mathcal N'\) and \(\pi\),
        \(\abs{\dom\mathcal N'}=\abs{\dom\mathcal N}\), and
        \(\dom\mathcal N'\setminus\dom\mathcal M\) is mapped by \(\pi\)
        bijectively onto \(D\); so \(\mathcal N'\) is a proper extension of
        \(\mathcal M\) iff \(j\) has proper range.
  \item If \((\mathcal N'_1,\pi_1)\) and \((\mathcal N'_2,\pi_2)\) both
        satisfy clause~1, then \(\pi_2^{-1}\circ\pi_1\) is an isomorphism
        \(\mathcal N'_1\to\mathcal N'_2\) that fixes every element of
        \(\dom\mathcal M\).
\end{enumerate}
The construction of \(\mathcal N'\) uses no form of the Axiom of Choice.
\end{hbtheorem}

The domain of \(\mathcal N'\) in the proof is
\(\dom\mathcal M\cup\{\langle r,d\rangle:d\in D\}\), where
\(\langle r,d\rangle\) is the Kuratowski pair and \(r\) is a set chosen so
that none of the new pairs already lies in \(\dom\mathcal M\); the
interpretations are transported from \(\mathcal N\) along the bijection that
sends \(a\in\dom\mathcal M\) to \(j(a)\) and \(\langle r,d\rangle\) to \(d\).
The element \(r\) is obtained by Russell's diagonal construction from the set
of first coordinates of pairs in \(\dom\mathcal M\), so no choice is made.

\begin{hbcorollary}{thm:p02-w02-p02-033-upward-ls-inclusion}{proved=annex:p02-w02-p02-033-lsincl}{Upward L\"owenheim--Skolem in inclusion form}
Assume the Axiom of Choice.  Let \(\mathcal M\) be an infinite
\(\Lang\)-structure and
\[
 \kappa\geq
 \max\bigl(\abs{\Lang},\abs{\dom\mathcal M},\aleph_0\bigr).
\]
Then there is an \(\Lang\)-structure \(\mathcal N'\) with
\(\mathcal M\preccurlyeq\mathcal N'\),
\(\dom\mathcal M\ne\dom\mathcal N'\), and
\(\abs{\dom\mathcal N'}=\kappa\).  Equivalently, the canonical
\(\Lang(\dom\mathcal M)\)-expansion of \(\mathcal N'\), interpreting each
\(c_a\) as \(a\) itself, satisfies \(\operatorname{ElDiag}(\mathcal M)\).
\end{hbcorollary}

\begin{hbwarning}{def:p02-w02-p02-033-forms-differ}{How the embedding form and the inclusion form differ}
The two forms make the same existence claim up to isomorphism, but they are
different statements about different objects.
\begin{enumerate}
  \item In the embedding form the domain of \(\mathcal N\) need not contain
        a single element of \(\dom\mathcal M\); the structure
        \(\mathcal M\) reappears only through its image
        \(j[\mathcal M]\), which is an elementary substructure of
        \(\mathcal N\) isomorphic to \(\mathcal M\) by
        \cref{lem:embedding-image} and
        \cref{thm:p02-w02-p02-033-elementary-maps-embed}.  In the proof of
        \cref{thm:p02-upward-ls-elementary-extension}, \(\dom\mathcal N\)
        consists of elements of a model supplied by compactness, unrelated
        to \(\dom\mathcal M\).
  \item Even when \(\dom\mathcal M\subseteq\dom\mathcal N\), an elementary
        embedding with proper range does not make the inclusion map
        elementary or proper; \cref{def:p02-w02-p02-033-doubling-example}
        has \(\dom\mathcal M=\dom\mathcal N\).
  \item The inclusion form identifies each \(a\in\dom\mathcal M\) with an
        element of \(\dom\mathcal N'\) that is literally \(a\).  A formula with
        parameters from \(\dom\mathcal M\) can then be evaluated in
        \(\mathcal M\) and in \(\mathcal N'\) at the same tuple, which is
        what \(\mathcal M\preccurlyeq\mathcal N'\) asserts.
  \item Passing from inclusion to embedding is immediate.  Passing from
        embedding to inclusion requires the construction of
        \cref{thm:p02-w02-p02-033-copy-over-image}: new elements disjoint
        from \(\dom\mathcal M\) must be produced, and the result is
        determined by \((\mathcal N,j)\) only up to isomorphism fixing
        \(\dom\mathcal M\) pointwise.
  \item Both forms say more than the inherited upward theorem
        \cref{thm:upward-ls}, which yields only a structure elementarily
        equivalent to \(\mathcal M\) and no map from \(\mathcal M\) into it.
\end{enumerate}
\index{elementary extension!inclusion and embedding forms}
\end{hbwarning}

\begin{hblemma}{thm:p02-w02-p02-033-pure-sets}{proved=annex:p02-w02-p02-033-pure}{Elementarity in the empty signature}
Let \(\Lang=\varnothing\), so that a structure is a nonempty set and the only
atomic formulas are equalities.
\begin{enumerate}
  \item If \(\dom\mathcal M\) is infinite and
        \(\dom\mathcal M\subseteq\dom\mathcal N\), then
        \(\mathcal M\preccurlyeq\mathcal N\).
  \item If \(\dom\mathcal M\) is infinite, every injection
        \(j:\dom\mathcal M\to\dom\mathcal N\) is an elementary embedding
        \(\mathcal M\to\mathcal N\).
\end{enumerate}
\end{hblemma}

\begin{hbexample}{def:p02-w02-p02-033-doubling-example}{An embedding with proper range that is not an inclusion}
Let \(\Lang=\varnothing\) and let \(\mathcal M=\mathcal N\) be the set
\(\N\) as a structure of the empty signature.  The map \(j(n)=2n\) is
injective, so by \cref{thm:p02-w02-p02-033-pure-sets} it is an elementary
embedding; its range, the even numbers, is proper.  Here
\(\dom\mathcal M=\dom\mathcal N\), yet the inclusion map is the identity and
has full range: the pair \((\mathcal N,j)\) is a proper extension in
embedding form, while \(\mathcal N\) is not a proper extension of
\(\mathcal M\) in inclusion form.

The construction of \cref{thm:p02-w02-p02-033-copy-over-image} repairs
this.  The set \(D\) is the set of odd numbers.  Let \(R\) be the set of
first coordinates of Kuratowski pairs \(\langle s,d\rangle\in\N\) with
\(d\in D\), and let \(r=\{x\in R:x\notin x\}\).  Then
\(\dom\mathcal N'=\N\cup\{\langle r,2m+1\rangle:m\in\N\}\), the two parts are
disjoint, and \(\pi(n)=2n\), \(\pi(\langle r,2m+1\rangle)=2m+1\) is a
bijection onto \(\N\).  The empty signature has nothing to transport, so
\(\mathcal N'\) is the set \(\dom\mathcal N'\) itself.  By
\cref{thm:p02-w02-p02-033-copy-over-image},
\(\mathcal M\preccurlyeq\mathcal N'\), the extension is proper, and
\(\abs{\dom\mathcal N'}=\aleph_0\); clause~1 of
\cref{thm:p02-w02-p02-033-pure-sets} confirms the elementarity directly.
This is the inclusion form of \cref{thm:p02-w02-p02-033-upward-ls-inclusion}
for \(\abs{\Lang}=0\) and \(\kappa=\aleph_0\).
\end{hbexample}

\begin{hbtheorem}{thm:p02-w02-p02-033-hypotheses-needed}{proved=annex:p02-w02-p02-033-hyp}{No hypothesis of upward L\"owenheim--Skolem can be dropped}
\begin{enumerate}
  \item If \(\mathcal M\) is finite, every elementary embedding
        \(j:\mathcal M\to\mathcal N\) is surjective.  Hence a finite
        structure has no proper elementary extension in either form.
  \item If \(\kappa<\abs{\dom\mathcal M}\), no structure of cardinality
        \(\kappa\) receives an elementary embedding from \(\mathcal M\).
  \item Let \(\Lang\) consist of constant symbols \(c_n\) for
        \(n\in\N\) and unary function symbols \(f_x\) for
        \(x\in\{0,1\}^{\N}\).  For a finite \(0\)--\(1\) sequence
        \(s=(s_0,\ldots,s_{\ell-1})\) put
        \(\operatorname{code}(s)=2^{\ell}+\sum_{i<\ell}s_i2^i\), and let
        \(\mathcal M\) have domain \(\N\), \(c_n^{\mathcal M}=n\) and
        \(f_x^{\mathcal M}(n)=\operatorname{code}(x(0),\ldots,x(n-1))\).
        Then \(\mathcal M\) is infinite and
        \(\abs{\dom\mathcal M}=\aleph_0<\abs{\Lang}\), and the domain of the
        target of every elementary embedding of \(\mathcal M\) with proper
        range has cardinality at least \(\abs{\mathcal P(\N)}\).  So for
        \(\kappa=\aleph_0\), which satisfies
        \(\kappa\geq\max(\abs{\dom\mathcal M},\aleph_0)\) but not
        \(\kappa\geq\abs{\Lang}\), \(\mathcal M\) has no proper elementary
        extension of cardinality \(\kappa\) in either form.
\end{enumerate}
\end{hbtheorem}

Clause~3 locates the use of \(\kappa\geq\abs{\Lang}\) in the proof of
\cref{thm:p02-upward-ls-elementary-extension}: downward
L\"owenheim--Skolem is applied in a signature of cardinality
\[
 \max\bigl(\abs{\Lang},\abs{\dom\mathcal M},\kappa\bigr),
\]
and its bound is \(\kappa\) only when \(\kappa\) dominates \(\abs{\Lang}\).  In the example,
every new element \(b\) is separated from each named number by a constant,
and the functions \(f_x\) then send \(b\) to \(\abs{\mathcal P(\N)}\)
pairwise distinct values.

\subsection{Nonstandard models by compactness}
\label{sec:p02-w02-p02-033-nonstandard}

Compactness builds a model in which an infinite family of conditions on new
constant symbols holds, provided each finite part of the family can be met in
a structure already at hand.  This subsection states the procedure, applies
it to three theories, and records where it fails.  In every example the
witness for an arbitrary finite fragment is written down and checked, and the
conclusions drawn from the resulting model are proved in the annex.

The arithmetic examples use the language
\(\Lang_A=\{0,S,+,\times\}\), its standard structure, numerals
\(\bar n=S^n0\), and the abbreviation \(r<t\) for
\(\exists z\,(r+Sz=t)\), all from \cref{def:language-arithmetic}
\cite[lines 8306--8438]{FrozenFormalLogicCompactnessApplications}.  Here that
structure is written \(\mathcal N_{\mathrm{ar}}=(\N,0,S,+,\times)\), as in
\cref{thm:p02-nonstandard-arithmetic-model}, where multiplication is written
\(\cdot\).  For a structure \(\mathcal M\) and a formula
\(\varphi(x_1,\ldots,x_k)\), \(\mathcal M\entails\varphi(b_1,\ldots,b_k)\)
means satisfaction under an assignment sending each \(x_i\) to \(b_i\); when
\(t_1,\ldots,t_k\) are closed terms, \(\varphi(t_1,\ldots,t_k)\) is the
sentence obtained by substituting them.  The models of
\(\operatorname{Th}(\mathcal N_{\mathrm{ar}})\), the set of
\(\Lang_A\)-sentences true in \(\mathcal N_{\mathrm{ar}}\), are classified as
standard or nonstandard by isomorphism type in \cref{def:nonstandard-model}
\cite[lines 9188--9261]{FrozenFormalLogicCompactnessApplications}.

\begin{hbmethod}{def:p02-w02-p02-033-compactness-method}{Building a model by compactness from finite witnesses}
\emph{Input}: an \(\Lang\)-structure \(\mathcal A\) and a set \(\Gamma\) of
sentences in \(\Lang\) enlarged by new constant symbols, expressing the
conditions the new elements are to meet.
\begin{enumerate}
  \item Choose the base theory.  When only sentences of \(\Lang\) are to be
        preserved, take \(B=\operatorname{Th}_{\Lang}(\mathcal A)\); for
        \(\mathcal A=\mathcal N_{\mathrm{ar}}\), where every element is the
        value of a numeral, this already yields an elementary embedding
        (clause~1 of \cref{thm:p02-w02-p02-033-numeral-map}).  When the
        conclusion must include an elementary embedding of an arbitrary
        \(\mathcal A\), take \(B=\operatorname{ElDiag}(\mathcal A)\) of
        \cref{def:p02-elementary-diagram}, which names every element \(a\) by
        its own constant \(c_a\).
  \item Form \(T=B\cup\Gamma\).
  \item For an \emph{arbitrary} finite \(\Delta\subseteq T\), interpret the
        finitely many new constants occurring in \(\Delta\cap\Gamma\) by
        elements of \(\mathcal A\) (and each \(c_a\) by \(a\)), and verify
        every sentence of \(\Delta\) in that expansion.  The interpretation
        may, and in the examples below must, depend on \(\Delta\).
  \item Assume the Axiom of Choice and apply compactness
        (\cref{thm:fol-compactness}, clause~1) to obtain a model
        \(\mathcal K\) of \(T\).
  \item Read off the conclusion in the reduct of \(\mathcal K\) to
        \(\Lang\): it satisfies \(\operatorname{Th}_{\Lang}(\mathcal A)\), and
        when \(B=\operatorname{ElDiag}(\mathcal A)\) the map
        \(a\mapsto c_a^{\mathcal K}\) is an elementary embedding by
        \cref{thm:diagram-tests}.  The new constants denote elements meeting
        every condition in \(\Gamma\).
  \item If wanted, apply downward L\"owenheim--Skolem
        (\cref{thm:downward-ls}) in the enlarged signature to bound the size,
        and \cref{thm:p02-w02-p02-033-copy-over-image} to pass to the
        inclusion form.
\end{enumerate}
\emph{Side conditions.}  Step~3 must cover every finite subset, not only an
increasing chain of them.  It fails when \(\mathcal A\) is finite and
\(\Gamma\) asks for an element different from every element of
\(\mathcal A\) (\cref{thm:p02-w02-p02-033-fs-fails}).  In the three examples
below the structure \(\mathcal A\) used in step~3 has no expansion that is a
model of \(T\), so the model of step~4 is a different structure.
\end{hbmethod}

\begin{hbdefinition}{def:p02-w02-p02-033-numeral-map}{Numeral map; standard and nonstandard elements}
Let \(\mathcal M\) be an \(\Lang_A\)-structure.  Its \emph{numeral map} is
\[
 \operatorname{num}_{\mathcal M}:\N\to\dom\mathcal M,\qquad
 \operatorname{num}_{\mathcal M}(n)\defeq\bar n^{\mathcal M}.
\]
An element of \(\dom\mathcal M\) is \emph{standard} when it lies in the range
of \(\operatorname{num}_{\mathcal M}\), and \emph{nonstandard} otherwise.
Write \(<^{\mathcal M}\) for the relation
\(\{(b,b'):\mathcal M\entails b<b'\}\) on \(\dom\mathcal M\).  For
\(n\in\N\) let
\[
 \theta_n\defeq\forall x\,\Bigl(x<\bar n\leftrightarrow
   \bigvee_{k<n}x=\bar k\Bigr),
 \qquad
 \tau_n\defeq\forall x\,\Bigl(\bigvee_{k\leq n}x=\bar k\lor\bar n<x\Bigr),
\]
where the empty disjunction (in \(\theta_0\)) is the sentence
\(\neg\,\bar 0=\bar 0\).
\index{nonstandard model!nonstandard element}
\end{hbdefinition}

\begin{hblemma}{thm:p02-w02-p02-033-numeral-map}{proved=annex:p02-w02-p02-033-numeral}{The numeral map of a model of true arithmetic}
Let \(\mathcal M\) be an \(\Lang_A\)-structure with
\(\mathcal M\entails\operatorname{Th}(\mathcal N_{\mathrm{ar}})\).
\begin{enumerate}
  \item \(\operatorname{num}_{\mathcal M}\) is an elementary embedding
        \(\mathcal N_{\mathrm{ar}}\to\mathcal M\).
  \item For \(b\in\dom\mathcal M\) and \(n\in\N\),
        \(\mathcal M\entails b<\bar n\) iff
        \(b=\operatorname{num}_{\mathcal M}(k)\) for some \(k<n\).  If \(b\)
        is nonstandard, then \(\mathcal M\entails\bar n<b\) for every
        \(n\in\N\).
  \item The following are equivalent: \(\mathcal M\cong\mathcal N_{\mathrm{ar}}\);
        \(\operatorname{num}_{\mathcal M}\) is surjective; \(\mathcal M\) has
        no nonstandard element.
  \item If \(b\) is nonstandard, there is exactly one \(b'\) with
        \(S^{\mathcal M}(b')=b\); it is nonstandard and \(b'<^{\mathcal M}b\).
        Hence, when \(\mathcal M\) has a nonstandard element, the nonempty set
        of nonstandard elements has no \(<^{\mathcal M}\)-least element,
        \(<^{\mathcal M}\) is not a well-ordering of \(\dom\mathcal M\), and
        recursion along the unique predecessor yields a sequence
        \((b_k)_{k\in\N}\) of nonstandard elements with
        \(b_{k+1}<^{\mathcal M}b_k\) for all \(k\).
  \item If \(\mathcal M\) has a nonstandard element, no \(\Lang_A\)-formula
        \(\varphi(x,y_1,\ldots,y_l)\) and parameters
        \(e_1,\ldots,e_l\in\dom\mathcal M\) satisfy
        \(\{b:\mathcal M\entails\varphi(b,e_1,\ldots,e_l)\}=
        \operatorname{num}_{\mathcal M}[\N]\).
\end{enumerate}
\end{hblemma}

Clause~3 turns the classification of \cref{def:nonstandard-model} into a test
on elements.  Clause~5 says that a model cannot single out its own standard
elements by a formula: the set of standard elements is an ambient set, visible
only from outside, as the external extension of
\cref{def:p02-w02-p02-033-internal-countability} is.

\begin{hbtheorem}{thm:p02-w02-p02-033-nonstandard-direct}{proved=annex:p02-w02-p02-033-direct}{A nonstandard model of true arithmetic from one new constant}
Assume the Axiom of Choice.  Let \(c\) be a new constant symbol and
\[
 T_1\defeq\operatorname{Th}(\mathcal N_{\mathrm{ar}})\cup
 \{\neg\,c=\bar n:n\in\N\}.
\]
\begin{enumerate}
  \item \emph{Finite satisfiability.}  For finite \(\Delta\subseteq T_1\),
        let \(m_\Delta=1+\max\{n:\neg\,c=\bar n\in\Delta\}\), with
        \(m_\Delta=0\) when no such sentence lies in \(\Delta\).  The
        expansion of \(\mathcal N_{\mathrm{ar}}\) interpreting \(c\) as
        \(m_\Delta\) satisfies \(\Delta\).
  \item \(T_1\) has a model whose domain is countable.
  \item Let \(\mathcal K\entails T_1\) and let \(\mathcal M\) be its
        \(\Lang_A\)-reduct.  Then
        \(\mathcal M\entails\operatorname{Th}(\mathcal N_{\mathrm{ar}})\),
        the element \(c^{\mathcal K}\) is nonstandard,
        \(\mathcal M\entails\bar n<c^{\mathcal K}\) for every \(n\),
        \(\operatorname{num}_{\mathcal M}\) is an elementary embedding with
        proper range, and \(\mathcal M\) is elementarily equivalent to
        \(\mathcal N_{\mathrm{ar}}\) but not isomorphic to it.
  \item There is an \(\Lang_A\)-structure \(\mathcal M'\) with countable
        domain, \(\mathcal N_{\mathrm{ar}}\preccurlyeq\mathcal M'\) and
        \(\N\ne\dom\mathcal M'\).
\end{enumerate}
\end{hbtheorem}

\Cref{thm:nonstandard-model} carries out the same finite-satisfiability step
for an arbitrary \(\Lang_A\)-theory satisfied by \(\mathcal N_{\mathrm{ar}}\);
clause~3 together with \cref{thm:p02-w02-p02-033-numeral-map} supplies, for
\(T=\operatorname{Th}(\mathcal N_{\mathrm{ar}})\), the agreement on sentences
and the failure of isomorphism set as \cref{ex:nonstandard-true-arithmetic}.
Clause~4 is the conclusion of \cref{thm:p02-nonstandard-arithmetic-model}; it
is reached here with one new constant, where that corollary adds
\(\aleph_0\) constants through
\cref{thm:p02-upward-ls-elementary-extension}.

The second example asks the new element to satisfy infinitely many conditions
of a different kind, so the finite witnesses are no longer large numbers.
For \(\Lang_A\)-terms \(t,u\), write \(t\mid u\) for
\(\exists z\,(t\times z=u)\), where \(z\) is the first variable occurring in
neither \(t\) nor \(u\).

\begin{hbtheorem}{thm:p02-w02-p02-033-divisible}{proved=annex:p02-w02-p02-033-divisible}{Divisibility by every prime}
Assume the Axiom of Choice.  Let \(c\) be a new constant symbol and
\[
 T_2\defeq\operatorname{Th}(\mathcal N_{\mathrm{ar}})\cup
 \{\neg\,c=\bar 0\}\cup\{\bar p\mid c:p\text{ a prime number}\}.
\]
\begin{enumerate}
  \item \emph{Finite satisfiability.}  For finite \(\Delta\subseteq T_2\),
        let \(m_\Delta\) be the product of the primes \(p\) with
        \(\bar p\mid c\in\Delta\), with \(m_\Delta=1\) when there are none.
        The expansion of \(\mathcal N_{\mathrm{ar}}\) interpreting \(c\) as
        \(m_\Delta\) satisfies \(\Delta\).
  \item \(T_2\) has a model.  If \(\mathcal K\entails T_2\) and
        \(\mathcal M\) is its \(\Lang_A\)-reduct, then
        \(b=c^{\mathcal K}\) satisfies \(\mathcal M\entails\neg\,b=\bar 0\)
        and \(\mathcal M\entails\bar p\mid b\) for every prime \(p\), while
        for each \(n\geq1\) there is a prime \(q\) with
        \(\mathcal M\entails\neg\,\bar q\mid\bar n\).  In particular \(b\) is
        nonstandard.
\end{enumerate}
\end{hbtheorem}

\begin{hbexample}{def:p02-w02-p02-033-fragment-witnesses}{Finite witnesses change with the fragment}
Let \(\sigma\) be any sentence of \(\operatorname{Th}(\mathcal N_{\mathrm{ar}})\).
\begin{enumerate}
  \item For \(\Delta=\{\sigma,\neg\,c=\bar 0,\neg\,c=\bar 3,\neg\,c=\bar 8\}
        \subseteq T_1\), the recipe of
        \cref{thm:p02-w02-p02-033-nonstandard-direct} gives
        \(m_\Delta=9\).  The fragment
        \(\Delta\cup\{\neg\,c=\bar 9\}\) is not satisfied by \(c\mapsto9\);
        its recipe value is \(10\).
  \item For \(\Delta=\{\sigma,\neg\,c=\bar 0,\bar 2\mid c,\bar 3\mid c,
        \bar 7\mid c\}\subseteq T_2\), the recipe of
        \cref{thm:p02-w02-p02-033-divisible} gives \(m_\Delta=42\), and
        \(\bar 7\mid c\) holds with \(z\mapsto6\).  The fragment
        \(\Delta\cup\{\bar 5\mid c\}\) is not satisfied by \(c\mapsto42\),
        since \(5\) does not divide \(42\); its recipe value is \(210\).
\end{enumerate}
No single natural number serves every fragment of \(T_1\), or every fragment
of \(T_2\) (\cref{thm:p02-w02-p02-033-fs-fails}).  Compactness needs only
one witness per fragment, and the nonstandard element of the resulting model
meets all the conditions at once, which no natural number does.
\end{hbexample}

The third example concerns a structure most of whose elements are not values
of closed terms, so the base theory is the elementary diagram.  Let
\(\Lang_{\mathrm{of}}=\{0,1,+,\cdot,<\}\), with constants \(0,1\), binary
function symbols \(+,\cdot\) and a binary relation symbol \(<\), and let
\(\mathcal R_{\mathrm{of}}=(\mathbb R,0,1,+,\cdot,<)\) be the ordered field of
real numbers.  For \(n\in\N\) let \(\underline n\) be the closed term with
\(\underline 0=0\) and \(\underline{n+1}=(\underline n+1)\).  An
\(\Lang_{\mathrm{of}}\)-structure \(\mathcal K\) is \emph{Archimedean} when
for every \(x\in\dom\mathcal K\) there is \(n\in\N\) with
\(\mathcal K\entails x<\underline n\).

\begin{hbtheorem}{thm:p02-w02-p02-033-infinitesimal}{proved=annex:p02-w02-p02-033-infinitesimal}{An elementary extension of the real field with infinitesimals}
Assume the Axiom of Choice.  Let \(c\) be a constant symbol not in
\(\Lang_{\mathrm{of}}(\mathbb R)\) and
\[
 T_3\defeq\operatorname{ElDiag}(\mathcal R_{\mathrm{of}})\cup
 \{c_a<c:a\in\mathbb R\}.
\]
\begin{enumerate}
  \item \emph{Finite satisfiability.}  For finite \(\Delta\subseteq T_3\),
        let \(r_\Delta=1+\max\{a:c_a<c\in\Delta\}\), with \(r_\Delta=0\)
        when there is no such sentence.  The expansion of
        \((\mathcal R_{\mathrm{of}})_{\mathbb R}\) interpreting \(c\) as
        \(r_\Delta\) satisfies \(\Delta\).
  \item \(T_3\) has a model.  If \(\mathcal K\entails T_3\) and
        \(\mathcal M\) is its \(\Lang_{\mathrm{of}}\)-reduct, then
        \(j(a)=c_a^{\mathcal K}\) defines an elementary embedding
        \(j:\mathcal R_{\mathrm{of}}\to\mathcal M\), and
        \(\mathcal M\entails j(a)<c^{\mathcal K}\) for every
        \(a\in\mathbb R\).  Moreover
        \(\mathcal M\entails c^{\mathcal K}<\underline n\) for no
        \(n\in\N\), so \(\mathcal M\) is not Archimedean.
  \item There is \(d\in\dom\mathcal M\setminus j[\mathbb R]\) with
        \(\mathcal M\entails0<d\) and \(\mathcal M\entails d<j(a)\) for every
        real \(a>0\).
  \item There are an \(\Lang_{\mathrm{of}}\)-structure \(\mathcal M'\) with
        \(\mathcal R_{\mathrm{of}}\preccurlyeq\mathcal M'\) and an element
        \(\varepsilon\in\dom\mathcal M'\setminus\mathbb R\) with
        \(\mathcal M'\entails0<\varepsilon\) and
        \(\mathcal M'\entails\varepsilon<a\) for every real \(a>0\).
\end{enumerate}
\end{hbtheorem}

The parameter-free theory would not serve here.  A model of
\(\operatorname{Th}_{\Lang_{\mathrm{of}}}(\mathcal R_{\mathrm{of}})\cup
\{\underline n<c:n\in\N\}\) is elementarily equivalent to
\(\mathcal R_{\mathrm{of}}\), but that theory names no irrational number, so
the construction would yield no map from \(\mathbb R\) into the model.  The
constants \(c_a\) of the elementary diagram supply the embedding, which is the
reason for the choice in step~1 of
\cref{def:p02-w02-p02-033-compactness-method}.

\begin{hblemma}{thm:p02-w02-p02-033-fs-fails}{proved=annex:p02-w02-p02-033-fails}{Where finite satisfiability fails}
\begin{enumerate}
  \item Let \(\mathcal F\) be an \(\Lang\)-structure with finite domain
        \(\{a_1,\ldots,a_k\}\) and let \(c\) be a new constant symbol.  The set
        \(\operatorname{ElDiag}(\mathcal F)\cup\{\neg\,c=c_a:a\in\dom\mathcal F\}\)
        has a finite subset with no model, and so has no model.
  \item Every finite subset of \(T_1\), and every finite subset of \(T_2\),
        is satisfied by an expansion of \(\mathcal N_{\mathrm{ar}}\), but no
        expansion of \(\mathcal N_{\mathrm{ar}}\) satisfies \(T_1\) or
        \(T_2\).  Every finite subset of \(T_3\) is satisfied by an expansion
        of \((\mathcal R_{\mathrm{of}})_{\mathbb R}\), but no such expansion
        satisfies \(T_3\).
\end{enumerate}
\end{hblemma}

Clause~1 is the compactness-side counterpart of clause~1 of
\cref{thm:p02-w02-p02-033-hypotheses-needed}: a finite structure has no proper
elementary extension, and the theory that would describe one is already
refuted by one finite fragment.  Clause~2 shows that compactness cannot be
replaced by a search for a single witness in the structure one started with.

\begin{hbcorollary}{thm:p02-w02-p02-033-not-axiomatizable}{proved=annex:p02-w02-p02-033-notax}{Three properties with no first-order axiomatization}
Assume the Axiom of Choice.
\begin{enumerate}
  \item No set \(\Sigma\) of \(\Lang_A\)-sentences has as its models exactly
        the \(\Lang_A\)-structures isomorphic to \(\mathcal N_{\mathrm{ar}}\).
  \item No set \(\Sigma\) of \(\Lang_A\)-sentences true in
        \(\mathcal N_{\mathrm{ar}}\) is such that \(<^{\mathcal M}\) is a
        well-ordering of \(\dom\mathcal M\) in every model \(\mathcal M\) of
        \(\Sigma\).
  \item No set \(\Sigma\) of \(\Lang_{\mathrm{of}}\)-sentences true in
        \(\mathcal R_{\mathrm{of}}\) has only Archimedean models.
\end{enumerate}
\end{hbcorollary}

Clause~2 has content because \(<^{\mathcal N_{\mathrm{ar}}}\) is the usual
order of \(\N\), which is a well-ordering; clause~3 has content because
\(\mathcal R_{\mathrm{of}}\) is itself Archimedean, every real number lying
below some natural number.

\begin{hbwarning}{def:p02-w02-p02-033-nonstandard-limits}{What the compactness construction does and does not provide}
\begin{enumerate}
  \item The construction proves existence only.  It does not describe the
        model of \(T_1\), \(T_2\) or \(T_3\), does not make it unique up to
        isomorphism, and gives no formula picking out the new element; by
        clause~5 of \cref{thm:p02-w02-p02-033-numeral-map}, a model of
        \(\operatorname{Th}(\mathcal N_{\mathrm{ar}})\) with a nonstandard
        element cannot even define its own standard elements.
  \item A nonstandard element is ``infinite'' only relative to the order
        \(<^{\mathcal M}\) of the model: it lies above every
        \(\operatorname{num}_{\mathcal M}(n)\).  In the countable model of
        clause~2 of \cref{thm:p02-w02-p02-033-nonstandard-direct} it is one
        element of a countable set.  Likewise the element \(d\) of
        \cref{thm:p02-w02-p02-033-infinitesimal} is positive and below every
        \(j(a)\) with \(a>0\) in the order of \(\mathcal M\); it is not a real
        number.
  \item Every model of \(T_1\) satisfies every sentence true in
        \(\mathcal N_{\mathrm{ar}}\), including each instance of the induction
        schema.  What fails in it is the induction principle of the ambient
        theory applied to the set \(\operatorname{num}_{\mathcal M}[\N]\),
        which is not one of the model's definable sets.
  \item Each existence clause above uses the Axiom of Choice of the ambient
        theory through \cref{thm:fol-compactness}.  The
        finite-satisfiability clauses,
        \cref{thm:p02-w02-p02-033-numeral-map} and
        \cref{thm:p02-w02-p02-033-fs-fails} do not.
\end{enumerate}
\index{nonstandard model!compactness construction}
\end{hbwarning}

\subsection{Internal and external countability}
\label{sec:p02-w02-p02-033-skolem}

The signature \(\Lang_{\in}\) of \cref{sec:zf-zfc-01} has the single binary
relation symbol \(\in\), so \(\abs{\Lang_{\in}}\) is finite.  For an
\(\Lang_{\in}\)-structure \(\mathcal M\) write \(\in^{\mathcal M}\) for the
interpretation of \(\in\).

\begin{hbdefinition}{def:p02-w02-p02-033-internal-countability}{External extension and the internal countability formulas}
For \(a\in\dom\mathcal M\), the \emph{external extension} of \(a\) in
\(\mathcal M\) is
\[
 \operatorname{ext}_{\mathcal M}(a)
 \defeq
 \{b\in\dom\mathcal M:\mathcal M\entails b\in a\}
 =\{b\in\dom\mathcal M:b\in^{\mathcal M}a\}.
\]
It is a set of the ambient set theory, a subset of \(\dom\mathcal M\); the
element \(a\) itself is an arbitrary object of the domain.  The following
\(\Lang_{\in}\)-formulas are abbreviations; \(\forall x\in a\,\psi\) and
\(\exists x\in a\,\psi\) abbreviate \(\forall x\,(x\in a\to\psi)\) and
\(\exists x\,(x\in a\land\psi)\).
\[
\begin{aligned}
 \mathrm{Emp}(e)&\defeq\forall t\,\neg(t\in e),\\
 \mathrm{Suc}(y,x)&\defeq\forall t\,\bigl(t\in y\leftrightarrow(t\in x\lor t=x)\bigr),\\
 \mathrm{Ind}(v)&\defeq\exists e\,(e\in v\land\mathrm{Emp}(e))\land
   \forall x\in v\,\exists y\in v\,\mathrm{Suc}(y,x),\\
 \mathrm{Om}(w)&\defeq\mathrm{Ind}(w)\land
   \forall v\,\bigl(\mathrm{Ind}(v)\to\forall x\,(x\in w\to x\in v)\bigr),\\
 \mathrm{Pw}(p,w)&\defeq\forall x\,\bigl(x\in p\leftrightarrow
   \forall t\,(t\in x\to t\in w)\bigr),\\
 \mathrm{OP}(z,u,v)&\defeq\forall y\,\Bigl(y\in z\leftrightarrow
   \bigl(\forall t\,(t\in y\leftrightarrow t=u)\\
 &\qquad\lor\forall t\,(t\in y\leftrightarrow(t=u\lor t=v))\bigr)\Bigr),\\
 \mathrm{In}(u,v,g)&\defeq\exists z\,\bigl(z\in g\land\mathrm{OP}(z,u,v)\bigr),\\
 \mathrm{Inj}(g,a,w)&\defeq
   \forall x\in a\,\exists n\in w\,\mathrm{In}(x,n,g)\\
 &\qquad\land\forall x\in a\,\forall n\in w\,\forall n'\in w\,
   \bigl(\mathrm{In}(x,n,g)\land\mathrm{In}(x,n',g)\to n=n'\bigr)\\
 &\qquad\land\forall x\in a\,\forall x'\in a\,\forall n\in w\,
   \bigl(\mathrm{In}(x,n,g)\land\mathrm{In}(x',n,g)\to x=x'\bigr),\\
 \mathrm{Ctbl}(a)&\defeq\exists w\,\bigl(\mathrm{Om}(w)\land\exists g\,\mathrm{Inj}(g,a,w)\bigr),\\
 \sigma_{\mathrm{unc}}&\defeq\exists w\,\exists p\,
   \bigl(\mathrm{Om}(w)\land\mathrm{Pw}(p,w)\\
 &\qquad\land\neg\exists g\,\mathrm{Inj}(g,p,w)\bigr).
\end{aligned}
\]
\(\mathrm{Om}(w)\) says that \(w\) is the least inductive set, the object
called \(\omega\) in \cref{sec:zf-zfc-01}; \(\mathrm{Pw}(p,w)\) says that
\(p\) is the power set of \(w\); \(\mathrm{In}(u,v,g)\) says
\(\langle u,v\rangle\in g\) for the Kuratowski pair; and
\(\mathrm{Ctbl}(a)\) says that some member of the model is an injection of
\(a\) into \(\omega\), the internal reading of the definition of a countable
set in \cref{sec:set-prerequisites-05}.  An element \(a\) is
\emph{internally uncountable} in \(\mathcal M\) when
\(\mathcal M\entails\neg\mathrm{Ctbl}(a)\), and \emph{externally countable}
when \(\operatorname{ext}_{\mathcal M}(a)\) is countable in the ambient set
theory.
\index{Skolem's paradox!external extension}
\end{hbdefinition}

The theorems below assume of an \(\Lang_{\in}\)-theory \(T\) only that
\(T\) contains extensionality, the power-set axiom, every instance of
separation, and the axiom of infinity, all as stated in
\cref{sec:zf-zfc-01}, with the abbreviations \(\varnothing\) and \(S(x)\) in
the axiom of infinity read through \(\mathrm{Emp}\) and \(\mathrm{Suc}\), so
that the axiom is the sentence \(\exists I\,\mathrm{Ind}(I)\).  The theories
\(\mathsf{ZF}\) and \(\mathsf{ZFC}\) of \cref{sec:zf-zfc-01} contain all of
these.

\begin{hbtheorem}{thm:p02-w02-p02-033-cantor-internal}{proved=annex:p02-w02-p02-033-cantor}{The theory proves that an uncountable set exists}
Let \(T\) be an \(\Lang_{\in}\)-theory containing the axioms just listed.
\begin{enumerate}
  \item Every model \(\mathcal M\) of \(T\) has exactly one element
        \(\omega^{\mathcal M}\) with
        \(\mathcal M\entails\mathrm{Om}(\omega^{\mathcal M})\) and exactly
        one element \(\mathcal P(\omega)^{\mathcal M}\) with
        \(\mathcal M\entails
        \mathrm{Pw}(\mathcal P(\omega)^{\mathcal M},\omega^{\mathcal M})\).
  \item Every model of \(T\) satisfies \(\sigma_{\mathrm{unc}}\), so
        \(T\entails\sigma_{\mathrm{unc}}\).
  \item Assume the Axiom of Choice.  Then
        \(T\derives_{\mathsf{ND}_{\mathsf{FOL}}}\sigma_{\mathrm{unc}}\), and
        by soundness every model of \(T\) satisfies
        \(\sigma_{\mathrm{unc}}\) because \(T\) proves it.  In particular
        every model \(\mathcal M\) of \(T\) satisfies
        \(\exists a\,\neg\mathrm{Ctbl}(a)\), witnessed by
        \(\mathcal P(\omega)^{\mathcal M}\).
\end{enumerate}
\end{hbtheorem}

Clause~2 is Cantor's diagonal argument from \cref{sec:set-prerequisites-05}
carried out inside an arbitrary model: from a putative injection
\(g\in\dom\mathcal M\) of \(\mathcal P(\omega)^{\mathcal M}\) into
\(\omega^{\mathcal M}\), separation produces the member of the model that
collects those \(n\) lying outside the subset that \(g\) sends to \(n\), and
that member defeats \(g\).  Every step uses an axiom of \(T\) as interpreted
in \(\mathcal M\), which is why the conclusion holds in every model,
whatever its size.

\begin{hbtheorem}{thm:p02-w02-p02-033-skolem-resolution}{proved=annex:p02-w02-p02-033-skolem}{Skolem's paradox: internally uncountable, externally countable}
Assume the Axiom of Choice.  Let \(T\) be as in
\cref{thm:p02-w02-p02-033-cantor-internal} and suppose that \(T\) has a
model.
\begin{enumerate}
  \item \(T\) has a model \(\mathcal M\) whose domain is countable.
  \item Let \(\mathcal M\) be any model of \(T\) with countable domain, and
        write \(w=\omega^{\mathcal M}\) and
        \(p=\mathcal P(\omega)^{\mathcal M}\).  Then
        \(\mathcal M\entails\neg\mathrm{Ctbl}(p)\), while
        \(\operatorname{ext}_{\mathcal M}(p)\subseteq\dom\mathcal M\) is
        countable.
  \item For every \(g\in\dom\mathcal M\), let
        \(R_g=\{(b,c):\mathcal M\entails\mathrm{In}(b,c,g)\}\).  Then
        \(\mathcal M\entails\mathrm{Inj}(g,p,w)\) iff
        \(R_g\cap\bigl(\operatorname{ext}_{\mathcal M}(p)\times
        \operatorname{ext}_{\mathcal M}(w)\bigr)\) is an injective function
        from \(\operatorname{ext}_{\mathcal M}(p)\) into
        \(\operatorname{ext}_{\mathcal M}(w)\).  Consequently no
        \(g\in\dom\mathcal M\) codes such an injection.
  \item Nevertheless there is, in the ambient set theory, an injective
        function from \(\operatorname{ext}_{\mathcal M}(p)\) into
        \(\operatorname{ext}_{\mathcal M}(w)\).
\end{enumerate}
\end{hbtheorem}

\begin{hbwarning}{def:p02-w02-p02-033-internal-external}{What ``uncountable'' means inside a model}
The model of \cref{thm:p02-w02-p02-033-skolem-resolution} satisfies
\(\neg\mathrm{Ctbl}(p)\) not by accident of its construction but because
\(T\) proves \(\sigma_{\mathrm{unc}}\)
(\cref{thm:p02-w02-p02-033-cantor-internal}).  That sentence quantifies over
the elements of the model only: it says that no \emph{element} \(g\) of
\(\dom\mathcal M\) has an \(R_g\) that injects the members of \(p\) into the
members of \(w\).  Countability of \(\operatorname{ext}_{\mathcal M}(p)\)
is a statement of the ambient set theory, and the injection of clause~4
is a set of the ambient theory that equals no \(R_g\) with
\(g\in\dom\mathcal M\).  The two claims therefore concern different
collections of functions and do not conflict.  Three limits apply.
First, everything is conditional on \(T\) having a model; for
\(T=\mathsf{ZFC}\) the existence of a model is a hypothesis of the theorem,
and nothing in this section proves it.  Second, clause~3 of \cref{thm:p02-w02-p02-033-cantor-internal} and
clause~1 of \cref{thm:p02-w02-p02-033-skolem-resolution} use the Axiom of
Choice of the ambient theory through completeness and downward
L\"owenheim--Skolem; clause~2 of
\cref{thm:p02-w02-p02-033-cantor-internal} and clauses~2--4 of
\cref{thm:p02-w02-p02-033-skolem-resolution} do not.  Third,
\(\operatorname{ext}_{\mathcal M}(w)\) need not be in bijection with \(\N\)
through the numerals \(0,1,2,\ldots\) of the model, and
\(\in^{\mathcal M}\) need not be the membership relation of the ambient
theory; the argument uses neither.
\end{hbwarning}

The section relates to its neighbours as follows.
\Cref{thm:p02-w02-p02-033-copy-over-image} supplies the step used in the
proof of \cref{thm:p02-nonstandard-arithmetic-model}, where the target of an
elementary embedding of \(\mathcal N_{\mathrm{ar}}\) is replaced by a copy
containing \(\N\) itself.  By \cref{thm:diagram-tests}, an inclusion-form
extension is the same thing as a model of \(\operatorname{ElDiag}(\mathcal M)\)
in which each \(c_a\) denotes \(a\), so
\cref{thm:p02-w02-p02-033-upward-ls-inclusion} is the elementary-diagram
form of upward L\"owenheim--Skolem with the hypotheses
\(\kappa\geq\max(\abs{\Lang},\abs{\dom\mathcal M},\aleph_0)\) and
\(\mathcal M\) infinite, each of which is needed by
\cref{thm:p02-w02-p02-033-hypotheses-needed}.
\Cref{thm:p02-w02-p02-033-nonstandard-direct} obtains the conclusion of
\cref{thm:p02-nonstandard-arithmetic-model} directly by compactness, and with
\cref{thm:p02-w02-p02-033-numeral-map} it completes the analysis that
\cref{thm:nonstandard-model} and \cref{ex:nonstandard-true-arithmetic} begin
for arbitrary theories satisfied by \(\mathcal N_{\mathrm{ar}}\);
\cref{thm:p02-w02-p02-033-divisible,thm:p02-w02-p02-033-infinitesimal} apply
the same procedure to other theories, the second through the elementary
diagram of \cref{def:p02-elementary-diagram}.
\Cref{thm:p02-w02-p02-033-skolem-resolution} is the downward counterpart for
set theory, and it supplies the argument behind
\cref{def:p02-skolem-paradox}.
